\section{Introduction}
\label{ch:introduction}

\subsection{Institutional questions as a problem of useful assistance}

A student looking for information about enrollment, recognition of credits, or a university procedure rarely begins with the name of the document that contains the answer. The question usually describes a situation: a requirement that is unclear, a form that must be submitted, or a next step that cannot be inferred from the available instructions. Answering well requires connecting that situation to the right source, identifying any conditions that matter, and explaining what the student can do. A response can sound clear and confident while failing at any of these tasks.

The finki-hub chatbot provides the setting for this thesis. It is a conversational assistant for students at the Faculty of Computer Science and Engineering (FCSE), Ss. Cyril and Methodius University in Skopje. Its knowledge resources include short frequently asked questions (FAQs), longer institutional documents, and separate facilities for accessing structured information. Students can interact with the assistant through web and Discord interfaces. Their questions include Macedonian written in Cyrillic, Macedonian expressed through Latin transliteration, and English. The resulting problem combines heterogeneous evidence with varied ways of expressing an information need.

The difficulty is not simply finding text that mentions the same topic. A short FAQ may describe a familiar procedure without every exception. A regulation may contain an authoritative condition several sections away from the most easily retrieved passage. A question may concern the current academic year, while the available source belongs to an earlier period. A student may ask about an individual decision that public documents cannot establish. In these situations, useful assistance depends on recognizing what the available evidence permits the assistant to say.

This thesis therefore treats effectiveness as a property of the response to a question. The desired response is either a supported and sufficiently complete answer or appropriate guidance when a direct answer is unavailable or unsuitable. Guidance can include a necessary clarification, a request to verify a changing detail, or a referral to the responsible institutional service. It must identify a meaningful next step without implying that the assistant has accessed private records or completed an official transaction. A referral can help a student proceed even though the underlying request remains unresolved.

This distinction also sets the boundary of the study. The work examines the information and guidance produced by the assistant. It does not measure staff-time savings, avoided emails, or reductions in administrative workload. Those are different outcomes requiring different observations. Likewise, generating a complete response stream does not establish that the student received a correct answer, and recorded use does not establish satisfaction. The evaluation must examine what the response says and how its claims relate to the evidence available to it.

\subsection{Why retrieval and architecture matter}

Retrieval-augmented generation (RAG) provides external evidence to a language model when it constructs an answer~\cite{lewis2020rag}. In an institutional assistant, this makes source selection part of the answer process: the model receives passages chosen for the current question rather than relying only on information represented in its parameters. The arrangement creates opportunities to inspect the evidence, retain source links, and update institutional material independently of the generator. It also creates several places where an apparently successful request can still produce poor assistance.

A document must first be acquired and converted without losing relevant content. It must then be divided into retrievable units, represented in the index, and selected for the question. The generation prompt must preserve the useful conditions in those units and distinguish evidence from instructions. Finally, the answer must reach the user with a suitable explanation of sources and limitations. An error in an early stage can survive all later stages. Conversely, a well-retrieved passage can still be misread, summarized incompletely, or attached to the wrong claim.

The registered title, \emph{A Reference Architecture for Building Domain-Specific Retrieval-Augmented Conversational Assistants}, reflects this broader system concern. A reference architecture identifies responsibilities, interfaces, and points where an implementation may vary. In this thesis, source preparation, retrieval, generation, channel delivery, and observation are described separately so that a response can be traced through them. The architecture supports the effectiveness investigation by making the tested path and its excluded components explicit. It is not presented as experimentally superior to competing architectures.

One concrete design choice concerns the treatment of FAQs and document passages. Their length and structure differ, yet a single score threshold can determine whether either source type enters the generation context. Lowering the threshold for document passages could recover conditions missed by a stricter shared threshold. It could also admit loosely related material that makes the answer harder to construct. Similarly, retaining the best candidate when every candidate falls below the threshold could supply useful evidence or merely prevent the system from recognizing that it lacks support. Neither effect follows from the amount of retrieved text alone.

These choices are suitable for a controlled comparison because they operate at an identifiable boundary. The candidate pool and scores can be captured once, the selection rules can change, and the remaining prompt and generation settings can stay fixed. The resulting comparison asks whether a change in the evidence supplied to the generator changes assistance outcomes. It does not require replacing the whole application or attributing every response difference to a new model.

\subsection{Research questions}

The central research problem is to determine what the available evidence establishes about the chatbot's ability to answer real student questions and provide useful guidance. Three questions organize the investigation.

\begin{description}
  \item[RQ1.] How often do the evaluated responses receive supported-answer or useful-routing labels, and what shortcomings remain?
  \item[RQ2.] How do separate FAQ and document score thresholds and a best-candidate fallback affect these outcomes when generation settings remain fixed?
  \item[RQ3.] How do question selection, reference construction, and AI-assessor behavior limit conclusions about effectiveness?
\end{description}

RQ1 separates complete factual assistance from helpful next steps and from failure categories. The distinction matters because a combined positive total can conceal whether a system actually answers questions or predominantly directs users elsewhere. Incomplete answers, unsupported material claims, uncertainty, and execution or assessment failures remain visible rather than being absorbed into a single average score.

RQ2 examines two retrieval decisions in a paired design. The same question is processed under a shared threshold or separate source thresholds, with fallback disabled or enabled. Comparing matched outcomes reduces variation attributable to the question set, while the four combinations distinguish the threshold change from the fallback change. The hypothesis that a lower document threshold might help is tested through answer outcomes as well as context diagnostics. An increase in selected passages is an explanation of the intervention, not the criterion for declaring it effective.

RQ3 treats assessment quality as part of the research problem. Automated judgments make it possible to inspect many responses, but they also introduce uncertainty about evidence interpretation, completeness, and category assignment. Repeated judgments and an additional assessor reveal sensitivity without supplying independent ground truth. The thesis asks what can still be concluded when those limitations are retained. The wording of these research questions was refined after the experiment to clarify the effectiveness focus; the frozen policies, hypotheses, outcomes, and labels were not changed.

\subsection{Study design and contribution}

The main evaluation contains 192 questions: 96 drawn from historical email groups and 96 from groups of questions selectively recorded by the deployed chatbot. A separate 20-case email pilot precedes the main analysis. Historical correspondence supplies realistic information needs, while chatbot questions expose shorter and sometimes context-dependent forms of interaction. The two sources are reported separately because equal sample sizes do not make their balanced mixture representative of actual use. Historical staff replies are not treated as authoritative expected answers.

The experiment reuses production retrieval and prompt code with controlled Luna generation. It captures evidence and prompts, then compares four frozen retrieval policies. Original questions are used without rewriting; conversation history, live provider-agent behavior, remote tools, and the browser or Discord delivery path are outside the controlled experiment. This boundary permits a focused comparison while preventing claims that the complete deployed conversational system has been evaluated. Architecture inspection and historical telemetry provide additional context, each with its own evidence limits.

The contribution has three connected parts. First, the thesis describes an inspectable architecture grounded in the chatbot implementation, including how sources become answer context and where structured access and delivery connect to that path. Second, it reports a controlled comparison of retrieval policies using real institutional questions and response-level outcomes. Third, it examines the reliability of the automated assessment through reference checks, repeated judgments, and cross-assessor comparison. These parts connect an implementation choice to student-facing behavior and then ask how strongly that behavior has been measured.

The completed evidence does not establish an improvement merely from admitting more document text. The shared-threshold configuration without fallback receives positive assistance labels on 31 of 96 email questions and 53 of 96 chatbot questions. Lowering the document threshold changes those totals to 30 and 52, respectively. The paired uncertainty intervals permit both improvement and deterioration. These are frozen AI-assessed outcomes, not validated factual accuracy rates. In particular, empty assessment references and disagreement between assessors limit the interpretation of absolute scores.

The reference architecture is correspondingly a grounded proposal for reuse, rather than a portability result from a second institution. The thesis does not claim a released public benchmark, completed language-model comparison, machine-translation experiment, or full conversational evaluation. Stating these boundaries is necessary to relate the completed work to the broader registered objectives without presenting planned extensions as results.

\subsection{Organization of the thesis}

The thesis proceeds from background and institutional data to architecture and implementation. Methodology defines the frozen retrieval-policy and prompt comparisons before results describe assistance outcomes, retrieval changes, the prompt comparison and assessment sensitivity. Discussion and conclusion answer the research questions and identify future work. The appendix explains retained evidence and reproducibility without including private correspondence or conversation records.
